\unitopen{26}
\begin{candos}
\item dobrać zaimek względny: rodzaj i liczba z poprzednika, przypadek z roli w zdaniu względnym;
\item opisać osobę, narzędzie lub system zdaniem względnym;
\item użyć zdania względnego z przyimkiem (\de{mit dem}, \de{an der}, \de{auf den});
\item rozumieć i tworzyć zdania względne w specyfikacjach.
\end{candos}

\section{Słownictwo}
\vocabtable{26}

\section{Zaimki względne}
Zaimek względny wygląda jak rodzajnik określony, z wyjątkiem celownika mnogiego (\de{denen}) i dopełniacza (\de{dessen, deren}, jednostka~27).

\begin{center}\small
\begin{tabular}{@{}lllll@{}}
\toprule
 & \textsf{\textbf{m}} & \textsf{\textbf{f}} & \textsf{\textbf{n}} & \textsf{\textbf{mnoga}}\\
\midrule
Nominativ & \de{der} & \de{die} & \de{das} & \de{die}\\
Akkusativ & \de{den} & \de{die} & \de{das} & \de{die}\\
Dativ & \de{dem} & \de{der} & \de{dem} & \de{denen}\\
\bottomrule
\end{tabular}
\end{center}
\textbf{Rodzaj i liczba} bierzesz z \emph{poprzednika}, \textbf{przypadek} z \emph{roli w zdaniu względnym}. Zdanie względne jest oddzielone
przecinkami, czasownik na końcu.

\begin{center}\small
\begin{tabular}{@{}p{0.34\linewidth}p{0.58\linewidth}@{}}
\toprule
\textsf{\textbf{Rola}} & \textsf{\textbf{Przykład}}\\
\midrule
podmiot (Nom) & \de{Das ist der Kollege, \textbf{der} das Projekt leitet.}\\
dopełnienie (Akk) & \de{Das ist die Lösung, \textbf{die} wir gesucht haben.}\\
dopełnienie (Dat) & \de{Das ist der Kunde, \textbf{dem} ich geholfen habe.}\\
z przyimkiem & \de{Das ist das Tool, \textbf{mit dem} wir arbeiten.}\\
 & \de{Das ist die Datenbank, \textbf{in der} die Daten liegen.}\\
\bottomrule
\end{tabular}
\end{center}

\begin{falle}
Polskie „który” zmienia się według roli i rodzaju, więc intuicja jest dobra, ale po niemiecku trzeba osobno ocenić
\emph{dwie} rzeczy. Typowy błąd: \de{der Kunde, den ich geholfen habe} (\de{helfen} wymaga celownika: \de{dem}). Drugi: opuszczenie
przecinka przed zaimkiem. Trzeci: \de{was} zamiast \de{das} po \de{alles}, \de{etwas}, \de{nichts}: \de{alles, was ich weiß}.
\end{falle}

\begin{szwajcaria}
W szwajcarskich dokumentach i umowach zdania względne są długie i formalne; rozbijaj je przy czytaniu na części. W mowie
używa się prostszych konstrukcji z \de{wo}: dialektalne \de{De Maa, wo i gseh han} odpowiada \de{der Mann, den ich gesehen habe}.
Rozpoznanie tego wzorca pomoże Ci w B1--B2 w słuchaniu.
\end{szwajcaria}

\begin{praca}
\de{Das ist der Service, der die Anfragen verarbeitet.} \de{Wir nutzen eine Plattform, die skaliert.} \de{Das Team, mit dem
ich zusammenarbeite, sitzt in Zürich.} \de{Die Funktion, die du suchst, ist in der neuen Version.} W opisie PR zdania względne
skracają listy.
\end{praca}

\section{Dialogi}
\textbf{Dialog 1 --- opis systemu}
\begin{dialog}
\turn{A}{Welche Anwendung meinst du?}{O jaką aplikację ci chodzi?}
\turn{B}{Die Anwendung, die die Rechnungen erstellt. Sie läuft auf dem Server, den wir letzte Woche gekauft haben.}{O aplikację, która tworzy faktury. Działa na serwerze, który kupiliśmy w zeszłym tygodniu.}
\turn{A}{Und wer ist der Entwickler, dem ich eine Frage stellen kann?}{A kto jest programistą, któremu mogę zadać pytanie?}
\turn{B}{Das ist Anna, mit der ich oft arbeite.}{To Anna, z którą często pracuję.}
\end{dialog}

\section{Ćwiczenia}
\begin{cwiczenia}
\exitem{Uzupełnij: \de{Das ist der Kollege, \dots{} das Projekt leitet.}}{\de{der}}
\exitem{Uzupełnij: \de{Das ist die Lösung, \dots{} wir gesucht haben.}}{\de{die}}
\exitem{Uzupełnij: \de{Das ist der Server, \dots{} wir gekauft haben.}}{\de{den}}
\exitem{Uzupełnij: \de{Das ist der Kunde, \dots{} ich geholfen habe.}}{\de{dem}}
\exitem{Uzupełnij: \de{Das ist das Tool, mit \dots{} wir arbeiten.}}{\de{dem}}
\exitem{Uzupełnij: \de{Das ist die Datenbank, in \dots{} die Daten liegen.}}{\de{der}}
\exitem{Uzupełnij: \de{Das sind die Kollegen, mit \dots{} ich arbeite.}}{\de{denen}}
\exitem{Połącz w jedno zdanie: \de{Das ist ein Server. Er läuft stabil.}}{\de{Das ist ein Server, der stabil läuft.}}
\exitem{Połącz: \de{Ich suche die Datei. Du hast sie gestern geschickt.}}{\de{Ich suche die Datei, die du gestern geschickt hast.}}
\exitem{Popraw: \de{der Kunde, den ich geholfen habe}.}{\de{der Kunde, dem ich geholfen habe}}
\exitem{Uzupełnij: \de{alles, \dots{} ich weiß}.}{\de{was}}
\exitem{Przetłumacz: „To jest narzędzie, którego używamy codziennie.”}{\de{Das ist das Tool, das wir jeden Tag benutzen.}}
\end{cwiczenia}

\weekend{Opisz trzy swoje narzędzia, projekty lub osoby w pracy zdaniami względnymi (po dwa zdania każde). Poproś AI,
by sprawdziła tylko przypadek zaimka. Nagraj 60 sekund.}
